% Full original nonexpressibility construction, with corrected indices and scope.

\section{Finite future bases and the complete-order obstruction}
\label{sec:annex-p09-gpss-complete-order-complement}

\begin{annexitem}{annex:p9.hr.ordinal.game}{thm:p09-hr-ordinal-game}{The ordinal interval strategy and its logical consequence}
\begin{proof}
We use elementary finite Cantor polynomials, including zero.
For \(r\geq0\), each such polynomial has a unique division
\(\delta=\omega^r\xi+\eta\), with \(\eta<\omega^r\):
separate its terms of exponent at least \(r\) from the
lower terms. If \(r\geq1\), every tail after a point of
\(\omega^r\alpha\), with \(\alpha>0\), has order type
\(\omega^r\gamma\) for some \(\gamma>0\).
Indeed write the point in its block
\(\omega^r\zeta+\eta\), \(\eta<\omega^r\).
The remainder of that block has type \(\omega^r\).
For \(r=1\) this is the infinite tail of the natural
numbers. For larger \(r\), a tail consists of the
remaining part of one \(\omega^{r-1}\)-block followed
by infinitely many such blocks; induction gives
type \(\omega^r\). Appending the remaining complete
\(\omega^r\)-blocks gives the asserted nonzero multiple.
All tails and initial segments remain finite Cantor
polynomials in the present application.

We prove the winning strategy by induction on \(n\).
At zero rounds there is no condition to violate.
For \(n\geq1\), write the two orders as
\(A=\omega^n\alpha+\rho\) and
\(B=\omega^n\beta+\rho\), with \(\rho<\omega^n\).
Consider the first selected point in either order;
the argument is symmetric, so suppose it is in \(A\).

If it belongs to the final copy of \(\rho\), choose
the point with the same position in \(B\)'s copy.
The intervals to their right are isomorphic.
Their left intervals have types
\(\omega^n\alpha+\delta\) and
\(\omega^n\beta+\delta\), with \(\delta<\rho\).
Divide \(\delta=\omega^{n-1}\zeta+\eta\).
Both left intervals are nonzero multiples of
\(\omega^{n-1}\) followed by the same
\(\eta<\omega^{n-1}\). The induction hypothesis
therefore supplies their \((n-1)\)-round strategy.

If the selected point lies in the initial
\(\omega^n\alpha\) part, let its initial-segment
type be \(\delta=\omega^{n-1}\xi+\eta\),
where \(\eta<\omega^{n-1}\).
When \(\xi=0\), choose the point \(\eta\) in
the first \(\omega^n\)-block of \(B\);
the left intervals are identical.
When \(\xi>0\), choose
\(\omega^{n-1}+\eta\), also in that first
block. Their left intervals are
\(\omega^{n-1}\xi+\eta\) and
\(\omega^{n-1}+\eta\), so the induction
hypothesis applies.

In this second case the right intervals,
after removing the two chosen points,
have types \(\omega^n\gamma+\rho\)
and \(\omega^n\gamma'+\rho\) for
nonzero \(\gamma,\gamma'\), by the
tail calculation. Divide \(\rho\)
at exponent \(n-1\); again these
are nonzero \(\omega^{n-1}\)-multiples
followed by the same smaller remainder.
The induction hypothesis supplies the
right-interval strategy too.

For the remaining \(n-1\) rounds use
the left or right strategy according
to the side on which the new point
is chosen. A repeated choice of the
first point is answered by its already
chosen mate. Each subgame receives
at most \(n-1\) moves; a winning
strategy for that many moves also
works for any smaller number.
The two subgames preserve order
internally, and the fixed first
pair preserves the order across
them. Repeated points inside a
subgame are handled by its strategy.
Thus equality and strict order are
preserved on the whole selected
tuple. This completes the induction.

For completeness, the game's
quantifier-depth consequence follows
by induction on formulas with assigned
variables among a selected tuple.
At depth zero, the winning condition
preserves each atomic order, equality
and unary statement, and hence every
Boolean combination of them.
For an existential formula of
positive depth, choose its witness
as the next move. The response and
the remaining strategy give the
induction hypothesis for its matrix.
The same argument with the roles
reversed gives the converse.
Negation supplies universal formulas,
and conjunction preserves the result.
This proves agreement on every
sentence of depth at most the
available number of rounds.

Finally, for \(k,m\geq n\) and
\(k,m\geq1\), take
\(\alpha=\omega^{k-n}\),
\(\beta=\omega^{m-n}\), and
\(\rho=0\). At \(n=0\) the
zero-round conclusion is immediate.
For \(n>0\) the displayed ordinal
strategy proves \(O_k\equiv_n O_m\).
The proof supplies the ordinal
comparison used in the source's
Theorem~3.3 and game appendix
\cite[pp.~9, 13--14]{HirshfeldRabinovich2003Future}.
\end{proof}
\end{annexitem}

\begin{annexitem}{annex:p9.hr.ordinal.tail}{thm:p09-hr-ordinal-tail}{Tail isomorphisms and predicate-free modal induction}
\begin{proof}
The tail of \(\omega\) from any of its points is
isomorphic to \(\omega\) by subtracting that
point's finite index. Suppose the result known
for \(O_{k-1}\). A tail of \(O_k\), viewed as
an \(\omega\)-sequence of \(O_{k-1}\)-blocks,
begins with the tail of one block and continues
with infinitely many whole blocks. The first
tail is itself \(O_{k-1}\) by induction.
The resulting order is \(O_k\). This proves
the assertion for every positive \(k\);
the open tail after a point is included
by beginning at its immediate successor.

A future formula has the value given
by the induced tail at its current
point. Isomorphism of these tails
therefore makes its value constant
on \(O_k\). Temporal formulas built
from future modalities are future:
induction on the syntax shows that
each operand's truth set on the
tail is its original truth set
restricted to the tail, after which
the future truth table gives the
same value. Constants and Boolean
operations preserve that property.

For the cross-model claim use
structural induction on a
predicate-free \(\mathsf{FTL}_d\)
formula. Constants agree; Boolean
operations preserve agreement.
At a modality \(\mu(B_1,\ldots,B_r)\),
each operand is constant on both
orders and has the same constant
value by induction. Substitute
true or false for its unary
argument in the supplied truth
table. The resulting one-variable
order formula \(\theta(t)\)
has depth at most \(d\).
It is future and homogeneous.
Its value is therefore the
truth value of \(\forall t\,\theta(t)\),
whose depth is at most \(d+1\).
The ordinal comparison transfers
that value whenever \(k,m\geq d+1\),
and in particular under the
source's stated margin \(k,m>d+1\).
This handles every finite arity
and every level of modal nesting
\cite[Lemmas~3.1--3.4]{HirshfeldRabinovich2003Future}.
\end{proof}
\end{annexitem}

\begin{annexitem}{annex:p9.hr.limit.separator}{thm:p09-hr-limit-separator}{The corrected limit hierarchy and predicate restriction}
\begin{proof}
All variables quantified in \(\varphi_n(x)\)
are restricted by their displayed guards to
the tail at \(x\). In the recursion, \(v\geq x\)
and \(u>v\) imply \(u>x\). The truth
of the formula is consequently unchanged
by deleting points below \(x\). The same
fact holds after restricting the quantifiers
to \(P\).

At \(x=0\), \(L_1(0,t)\) says that \(t\)
is a nonzero limit point of its predecessors.
For a finite Cantor polynomial this is
equivalent to its coefficient of \(\omega^0\)
being zero and at least one higher coefficient
being nonzero. If its last coefficient is
positive, \(t\) has an immediate predecessor,
which falsifies the universal condition.
If the last coefficient is zero and
\(t>0\), points below \(t\) approach
it by increasing a lower coefficient;
there is a point between every earlier
point and \(t\). This proves the first
level and excludes zero explicitly.

Inductively suppose \(L_j(0,u)\)
holds exactly at positive multiples
of \(\omega^j\). The positive
multiples of \(\omega^j\) below
\(t\) are cofinal in \(t\)
exactly when \(t\) is a positive
multiple of \(\omega^{j+1}\).
To verify this from the polynomial,
if any term of exponent less
than \(j\) is nonzero, a final
interval below \(t\) has no
\(\omega^j\)-multiple. If those
lower terms vanish but the
coefficient at exponent \(j\)
is positive, there is again
a final interval with no such
multiple: the preceding
\(\omega^j\)-multiple, or zero
when there is none positive,
is a lower cutoff for that
interval. If all terms through
exponent \(j\) vanish and
\(t>0\), use the least nonzero
exponent above \(j\); increasing
the lower \(\omega^j\) coefficient
below its final block gives
positive \(\omega^j\)-multiples
beyond every point below \(t\).
These observations prove the
induction step directly.

It follows that no positive
point below \(\omega^k\)
satisfies \(L_n\) if \(n\geq k\).
If \(n<k\), the point
\(\omega^n<\omega^k\) does.
The existential in \(\varphi_n\)
therefore gives exactly \(k>n\).

Predicate restriction is proved
by induction on the original
first-order formula: atoms have
their induced-order meaning,
Boolean clauses commute, and
the restricted existential or
universal ranges exactly over
the induced domain \(P\).
With \(x\in P\), this gives
the stated interpretation of
\(\varphi_n^P\). This construction
keeps the fixed lower parameter
inside every limit operation;
it is a future formula on
arbitrary labelled chains,
not merely on the chosen
ordinal examples.
\end{proof}
\end{annexitem}

\begin{annexitem}{annex:p9.hr.continuous.models}{thm:p09-hr-continuous-models}{The real embedding, tail types, and all game coordinates}
\begin{proof}
We verify the order properties
before using an isomorphism.
The lexicographic order \(C_k\)
has a least point and no greatest
point. It is dense: points in
one block have another point
between their real coordinates;
between points in different
blocks, choose a larger coordinate
in the earlier block.
The countable set
\(\omega^k\times([0,1)\cap\mathbb Q)\)
is order dense.

The order is complete. For a
nonempty bounded-above set,
take the supremum \(\gamma\)
of its block indices in
\(\omega^k\); an upper bound
ensures that this ordinal
is below \(\omega^k\).
If there is no largest
occurring block index,
\((\gamma,0)\) is the
supremum. If there is a
largest index \(\gamma\),
take the real supremum
of the coordinates occurring
in that block. If it is
less than 1 the corresponding
point of that block is
the supremum; if it is
1, the supremum is
\((\gamma+1,0)\).
The latter block exists,
since \(\omega^k\) is a
limit ordinal. Its being
no higher than any proposed
upper bound follows from
the coordinates approaching
1. For an infimum, a
nonempty set of indices
has a least member; the
infimum of the coordinates
in that block lies in
\([0,1)\) and gives the
infimum of the whole set.
These descriptions also
verify leastness and
greatestness of the bounds.

Here is the required real
isomorphism argument.
Any two countable dense
orders with least point
and no greatest point
are isomorphic by back
and forth: match their
least points, enumerate
both sets, and alternately
extend the finite increasing
correspondence to the next
unmatched point on each
side. Density supplies an
unused point between its
finite lower and upper
constraints, and absence
of a greatest point handles
an empty upper constraint.
Apply this to the displayed
countable dense subset
and \(\mathbb Q_{\geq0}\).
Extend its isomorphism
to \(C_k\) by taking
the supremum in
\(\mathbb R_{\geq0}\)
of the images of dense
points below or equal
to a given point.
The set is nonempty
and bounded by the
image of a larger
dense point. Two
dense points between
distinct original
points show that
the extension is
strictly increasing.
For any nonnegative
real \(r\), the
preimages of rationals
below \(r\) have a
supremum in \(C_k\),
with zero treated as
the least point;
density gives image
exactly \(r\).
This proves surjectivity
and the claimed
order isomorphism.

The marked points
\((\gamma,0)\) have
exactly the order
of their indices.
At a marked point,
the remaining block
indices form a tail
of \(\omega^k\),
isomorphic to
\(\omega^k\) by
the preceding lemma;
copying block coordinates
gives the labelled
tail isomorphism.
At an unmarked point
\((\gamma,r)\), \(r>0\),
the initial part of
its tail is the
unmarked interval
\([r,1)\).
An increasing affine
map identifies any
two such intervals.
After that part,
the remaining marked
blocks again have
index type \(\omega^k\).
This proves both
tail types. Future
locality then makes
the truth of a
one-predicate future
formula constant on
each type. Two
Boolean values give
exactly the four
displayed descriptions.

To lift the comparison
strategy, suppose the
next chosen point is
\((\gamma,r)\).
Play \(\gamma\) in
the ordinal comparison
game and let \(\delta\)
be its response.
Answer with \((\delta,r)\).
Repeat this at every
round, whichever
structure supplied
the chosen point.
The ordinal strategy
preserves equality
and order of the
first coordinates,
including repetitions.
When first coordinates
are equal, the
copied real coordinates
preserve their internal
order and equality.
When they differ,
lexicographic order
is decided by the
first coordinates.
Finally a point
is marked exactly
when its copied
second coordinate
is zero. The
complete selected
tuples therefore
preserve all atomic
relations and \(P\).
The game-to-formula
induction already
proved gives the
claimed \(n\)-equivalence
\cite[Lemmas~4.2--4.3 and appendix~A]{HirshfeldRabinovich2003Future}.
\end{proof}
\end{annexitem}

\begin{annexitem}{annex:p9.hr.modal.invariance}{thm:p09-hr-modal-invariance}{The two-colour induction with the quantifier margin retained}
\begin{proof}
Every \(\mathsf{FTL}_d\)
formula is future by
the substitution induction
for future truth tables.
It therefore has one
of the four truth-set
descriptions on each
model \(C_k\).
We prove that the
description is uniform
for \(k>d\), by
structural induction.

The proposition \(P\)
has its named
description; constants
have theirs. Boolean
operations on the
two truth values
are pointwise and
independent of \(k\).
For a modality
\(\mu(B_1,\ldots,B_r)\),
the induction hypothesis
replaces each operand
by the same one of
\(\top,\bot,P,\neg P\)
on every model in
the indicated range.
Substitute these
quantifier-free
descriptions into
its supplied
truth table.
The resulting
\(\theta(t;P)\)
still has quantifier
depth at most \(d\)
and is future.

The value of
\(\theta\) on
marked points
is its uniform
value on that
nonempty point type,
and is characterized
by the sentence
\[
                \forall t\,(P(t)\to\theta(t;P)).
\]
The same statement
for unmarked points
uses \(\neg P(t)\).
Both types are
nonempty in every
\(C_k\), so these
sentences recover
the two truth
values exactly.
Their depth is
at most \(d+1\).
For \(k,m>d\),
the previous
comparison lemma
applies at \(d+1\),
since these integers
satisfy \(k,m\geq d+1\).
Both sentences
therefore agree
in \(C_k,C_m\).
The description
of the modality
result is the
same in every
model. This
completes the
induction.

The argument bounds
only the substituted
truth table and
its single outer
universal quantifier.
It never bounds
the depth of the
whole temporal
formula, and thus
includes arbitrarily
deep nesting and
the infinite basis
of \(\mathsf{FTL}_d\).
\end{proof}
\end{annexitem}

\begin{annexitem}{annex:p9.hr.no.finite.future.basis}{thm:p09-hr-no-finite-future-basis}{The full finite-basis contradiction and its domains}
\begin{proof}
The marked subset
of \(C_k\) is
isomorphic to
\(O_k\), and its
least point is
marked. Predicate
restriction and
the separator
lemma therefore
give
\[
       C_k,(0,0)\models\varphi_n^P
                      \quad\Longleftrightarrow\quad k>n.
\]
For a fixed \(d\),
take \(n=d+1\),
and compare
\(C_{d+1}\) with
\(C_{d+2}\).
The displayed
formula has
different values,
while the
modal-invariance
lemma says that
every \(\mathsf{FTL}_d\)
formula in \(P\)
has the same
value at their
least points.
No such temporal
formula can be
equivalent to
\(\varphi_{d+1}^P\).
If a proposed
formula also uses
extraneous proposition
letters, fix all
of them false.
Replacing them
syntactically by
false reduces it
to the same
one-predicate
argument. Thus
extra letters
cannot repair
an equivalence
required for
all interpretations.

Each \(C_k\)
is a labelled
nonnegative real
line by the
proved isomorphism.
For the full
real line, use
that interpretation
at and after
zero and fix
the predicates
below zero.
The future
formula and
every future
modality formula
retain their
values at zero.
This proves
the result
on both real
domains.

Every finite
set of
first-order
future modalities
has a finite
maximum \(d\)
of the depths
of its supplied
truth tables;
increase it to
1 if necessary.
Its temporal
language is
contained in
\(\mathsf{FTL}_d\).
The separating
future formula
just constructed
therefore proves
the claimed
failure of
expressive
completeness.

For the ordinal
version use
the constant-only
induction and
choose exponents
\(d+2,d+3\).
Both exceed
\(d+1\), so
all constant-only
\(\mathsf{FTL}_d\)
formulas agree.
The future
order formula
\(\varphi_{d+2}\)
distinguishes
the two ordinals.
Any auxiliary
predicate letters
are again fixed
false. A finite
future basis
therefore cannot
be complete
over the class
of ordinals.

For the rational
version replace
each block
\([0,1)\) by
\([0,1)\cap\mathbb Q\).
The product
is countable,
dense, has a
least point
and no greatest
point. The
countable back
and forth
already given
identifies it
with \(\mathbb Q_{\geq0}\);
no completeness
argument is
needed here.
The marked
index order,
the two tail
types, and
the game that
copies the
second coordinate
are unchanged.
The modal
induction and
the restricted
ordinal separator
therefore give
the same
contradiction.
Prepending the
negative rational
half-line gives
the full
rational result
by future
locality. This
supplies the
source's
Theorem~5.1
with its
explicit
domain change.
\end{proof}
\end{annexitem}

\begin{annexitem}{annex:p9.gpss.until.complete.order.obstruction}{thm:p09-gpss-until-complete-order-obstruction}{The strict-until counterinstance and the impossible complement closure}
\begin{proof}
The strict-until truth
table has one
existential witness
and one nested
universal gap
quantifier, hence
depth 2.
The preceding
theorem at \(d=2\)
makes \(\varphi_3^P\)
an explicit
counterinstance,
using \(C_3,C_4\).
Both underlying
orders are complete.
The first-order
formula uses only
one distinguished
free variable
and quantifiers
in its future,
so it is in
the exact GPSS
input fragment.

There is also the shorter
counterinstance in the
body's worked example.
In \(C_k\), every open
interval between distinct
points contains an
unmarked point: choose
a positive coordinate
inside the earlier
block, or between the
two coordinates when
they share a block.
Thus an interval
operand false or \(P\)
cannot hold throughout
any strict-until gap.
An empty witness set
also makes until false.
For interval operand
true, every nonempty
one of the four witness
sets has a future
point, since marked
and unmarked points
both occur arbitrarily
far in the future.
For interval operand
\(\neg P\), at any
\((\gamma,r)\) the
next marked point
\((\gamma+1,0)\)
exists and the open
gap to it is entirely
unmarked. It supplies
a \(P\)-witness; a
point still inside
the current block
supplies a \(\neg P\)
or true witness.
These verify every
entry of the table.

The four truth sets
are Boolean closed,
and the table is
the same for every
\(k\geq1\).
An induction on
until-only formulas
therefore preserves
the same description
across all these
models, a sharper
conclusion for this
particular modality
than the general
depth argument.
In \(P_1\), which
has type \(\omega\),
every future marked
point has a preceding
marked point at
or above the initial
point. Using that
point as \(v\)
falsifies the
displayed formula.
In \(P_2\), the
point indexed by
\(\omega\) has
marked points
strictly between
every earlier
marked point
and itself.
It witnesses the
formula. This
proves the
worked rank-three
counterinstance
directly.

Suppose now that
all relative
complements of
left-labelled
bounded decompositions
had finite
left-labelled
decomposition
presentations
on complete
orders. Finite
conjunction and
existential
concatenation
are already
proved for
that class.
The empty
conjunction
and disjunction
are represented
by true and
false as
before.

Induct on
the exact
GPSS regional
stages. A
positive guarded
existential
conjoins its
finitely many
left and
right regional
children and
point literals.
By induction
their finite
lists are
available.
Positive
conjunction,
distribution
and concatenation
therefore give
a list for
that existential.
For its negative
sign, complement
each disjunct
and conjoin
the finite
lists. On
bounded intervals
this uses
the assumed
closure. On
future rays
it uses the
proved ray
complement
construction.
Every recursive
regional argument
has an actual
lower endpoint,
including the
quantified
midpoint.
Thus no formal
left endpoint
is used.

This induction
would represent
every \(\operatorname{Aft}_G(x)\)
component. The
proved future
special normal
form leaves
only those
components and
present literals.
Translate each
ray list to
\(L(U)\) and
retain its
left label,
then combine
the finite
Boolean operations.
The result
would express
every future
first-order
formula on
complete orders
in \(L(U)\),
contradicting
\(\varphi_3^P\).
Hence the
assumed bounded
closure is
false.

Even assuming
that closure
only for
original
endpoint-free
words, with
outputs permitted
in the enlarged
class, would
be enough:
the complement
of \(\alpha(x)\land D(x,y)\)
under \(x<y\)
is the disjunction
of \(\neg\alpha(x)\)
on the true
gap and
the assumed
complement of
\(D\). Thus
some original
word's bounded
complement also
fails in the
enlarged output
class. The
positive-chain
case already
proved is
consistent with
this conclusion;
it imposes no
general final
gap.
\end{proof}
\end{annexitem}

\begin{annexitem}{annex:p9.hr.guarded.future}{thm:p09-hr-guarded-future}{The exact protection provided by a strict future guard}
\begin{proof}
Fix a model
\(\mathcal M\)
and an actual
point \(x\).
Let \(\mathcal N\)
be its induced
tail at \(x\).
We prove by
structural
induction that
every formula
in \(\mathcal A\)
agrees in
the two models
at every
\(t>x\).
Atoms and
Boolean
operations
are immediate.
A strict
until at
\(t>x\)
uses only
witnesses and
gap points
strictly above
\(t\), where
the induction
hypothesis
applies.

Consider
\(K^-B(t)\).
If it holds
in \(\mathcal M\),
then for every
\(r\) in
\(\mathcal N\)
below \(t\),
its witness
lies strictly
above \(r\geq x\).
The induction
hypothesis
for \(B\)
transfers that
witness to
\(\mathcal N\).
Conversely
suppose it
holds in
\(\mathcal N\).
For a cut
\(r\geq x\)
use its
witness there
and the
induction
hypothesis.
For a cut
\(r<x\)
use instead
the instance
at the
actual point
\(x<t\)
in \(\mathcal N\).
Its witness
\(u\) satisfies
\(r<x<u<t\)
and has the
same \(B\)
truth in
both models.
This proves
the original
universal
condition for
every earlier
cut. Adjacency
causes no
exception:
if \(x,t\)
are adjacent,
that instance
at \(x\)
makes both
relevant
conditions
false.

At the
initial point
\(x\), every
strict-until
formula with
arguments in
\(\mathcal A\)
uses argument
evaluations
strictly after
\(x\), where
the result
just proved
applies.
Present atoms
are unchanged,
and Boolean
closure preserves
agreement.
Thus every
guarded future
formula agrees
with its tail
restriction
at \(x\).

The unguarded
past limit
does not have
that property.
On \(\mathbb R\)
at zero,
interpret
\(P=(-1,0)\)
in one model
and \(P=\varnothing\)
in the other.
They agree
at and after
zero. Every
left neighborhood
of zero has
a \(P\)-point
in the first
model, whereas
none does in
the second.
Thus \(K^-P\)
distinguishes
them. The
source's guarded
fragment may
contain such
a subformula
under until;
the subformula
itself is
not thereby
a future
formula.
\end{proof}
\end{annexitem}
